\chapter{形容词}
\section{形容词变格}

\subsection{一个原理}
\begin{center}
    \begin{myquote}
        做表语时用原型，做定语时才变格。
        限定名词性数格,最小冗余补冠词。
    \end{myquote}
\end{center}

\subsection{三种变格}
\begin{itemize}
    \item 强变：零冠词，形担任定冠词角色，变格同\hyperref[sec:pron:strong_adj]{定} \\
    \item 弱变：定冠词后，形仅做额外补充，变格见\hyperref[sec:pron:seak_adj]{表} \\
    \item 混合变：不定冠词后，形做最小冗余补充 \\
          不定冠词5种情况有歧义，形强变\\
          其他情况，形弱变
\end{itemize}
例外：二格强变不是-es而是-en

\subsection{其他代冠}
\begin{itemize}
    \item 相当于零冠词的代词 \\
          (只接复数或不可数，都是零冠词的情况，无法体现名词的格、性，即没有定冠词的“定性”)
          \begin{itemize}
              \item 不定代词：一些，只接复数，强调数量
                    \begin{itemize}
                        \item einige (a few)
                    \end{itemize}
              \item 不定代词：多少，
                    \begin{itemize}
                        \item wenige (few) + 复数
                        \item viele (many) + 复数
                        \item wenig (little) + 不可数
                        \item viel (much) + 不可数
                    \end{itemize}
          \end{itemize}
    \item 相当于定冠词代词
          \begin{itemize}
              \item 分配代词
                    \begin{itemize}
                        \item all- (all)
                        \item jed- (every)
                        \item beid- (both)
                    \end{itemize}
              \item 不定代词：某些
                    \begin{itemize}
                        \item manch- (some) ，可接单复，强调存在 \\
                              (some as in sometime not some time)
                    \end{itemize}
              \item 指示代词
                    \begin{itemize}
                        \item dies- (this)
                        \item solch- (such)
                    \end{itemize}
          \end{itemize}
    \item 相当于不定冠词的代词
          \begin{itemize}
              \item 否定冠词: kein
              \item 物主冠词: mein/dein/sein/ihr/unser/euer
          \end{itemize}
\end{itemize}

\subsection{特殊情况}
\begin{center}
    \begin{tabular}{|c|c|}
        \hline
        \rowcolor{gray}
        adj      & declension \\
        \hline
        -el, -er & 去e加词尾      \\
        \hline
        hoch     & 去c加词尾      \\
        \hline
        -a       & 无需加词尾      \\
        \hline
        城市名      & -er做结尾     \\
        \hline
    \end{tabular}
\end{center}